\documentclass[10pt]{extarticle}
\usepackage[a4paper,margin=1.5cm]{geometry}
\usepackage[T1]{fontenc}
\usepackage{amsmath,amssymb,mathtools}
\usepackage{graphicx}
\usepackage[export]{adjustbox}
\usepackage{multicol}
\usepackage{enumitem}
\usepackage{booktabs,array,tabularx}
\setlength{\columnsep}{0.7cm}
\setlength{\parindent}{0pt}
\setlength{\parskip}{3pt}
\pagestyle{empty}
\makeatletter
\newcommand{\qfigdef}[4]{\@namedef{qf@file@#1}{#2}\@namedef{qf@w@#1}{#3}\@namedef{qf@wide@#1}{#4}}
\newcommand{\qf@check}[1]{\@ifundefined{qf@file@#1}{\PackageError{score2}{Unknown figure id #1}{}}{}}
\newcommand{\qf@img}[4][]{\includegraphics[width=#2\dimexpr\@nameuse{qf@w@#4}pt\relax,max width=#3,max height=0.45\textheight,#1]{\@nameuse{qf@file@#4}}}
\newif\ifqf@full
\newif\ifqf@boxed
\newsavebox{\qf@box}
\newcommand{\qf@open}{\global\qf@fulltrue\par\global\qf@boxedtrue
  \setbox\qf@box\vbox\bgroup\hsize\textwidth\linewidth\textwidth}
\newcommand{\qf@emit}[1]{\par\egroup\global\qf@boxedfalse
  \ifdim\dimexpr\ht\qf@box+\dp\qf@box\relax<#1\textheight\noindent\usebox{\qf@box}\par\else\unvbox\qf@box\fi}
\newcommand{\qfullsplit}{\ifqf@boxed\qf@emit{0.85}\fi}
\newcommand{\qf@close}{\ifqf@boxed\qf@emit{0.45}\fi\global\qf@fullfalse}
\newcommand{\qfullbegin}{\end{multicols}\qf@open}
\newcommand{\qfullend}{\qf@close\begin{multicols}{2}}
\newcommand{\qfullstart}{\qf@open}
\newcommand{\qfulldone}{\qf@close}
\newcommand{\qfig}[2][1]{\qf@check{#2}%
  \ifqf@full\qf@blk{#1}{#2}\else
  \ifnum\@nameuse{qf@wide@#2}=1 \expandafter\@firstoftwo\else\expandafter\@secondoftwo\fi
  {\end{multicols}\begin{center}\qf@img{#1}{\textwidth}{#2}\end{center}\begin{multicols}{2}}%
  {\qf@blk{#1}{#2}}\fi}
\newcommand{\qf@blk}[2]{\par\smallskip{\centering\qf@img{#1}{\linewidth}{#2}\par}\smallskip}
\newcommand{\qfigi}[2][1]{\qf@check{#2}\qf@img[valign=c]{#1}{\linewidth}{#2}}
\makeatother
\newcommand{\papersec}[1]{\par\bigskip{\centering\Large\bfseries PAPER-#1\par}\medskip}
\newcommand{\partsec}[2]{\par\medskip\noindent{\large\bfseries PART-#1: #2}\par\smallskip}
\newcommand{\sect}[1]{\par\smallskip\noindent\textbf{#1}\par}
\newcommand{\qn}[1]{\par\smallskip\noindent\textbf{#1.}~}
\newcommand{\paraq}[2]{\par\smallskip\noindent\textbf{Paragraph for Questions #1 and #2}\par}
\newenvironment{qoptions}{\begin{itemize}[leftmargin=1.8em,itemsep=1pt,topsep=2pt,label={}]}{\end{itemize}}
\newcommand{\opt}[2]{\item[(#1)] #2}
\newcolumntype{L}[1]{>{\raggedright\arraybackslash}p{#1}}
\newcommand{\fwstart}{\end{multicols}\par\noindent\begin{minipage}{\textwidth}\small\setlength{\tabcolsep}{4pt}\renewcommand{\arraystretch}{1.5}}
\newcommand{\fwend}{\end{minipage}\par\medskip\begin{multicols}{2}}
% figure registry (generated by score2_tools.py finalize)
\qfigdef{p1_phy_s1_q01_1}{Pre-Score (2023-24) - Test 07 (29 Nov 2023)-img/p1_phy_s1_q01_1.png}{105.6}{0}
\qfigdef{p1_phy_s1_q02_1}{Pre-Score (2023-24) - Test 07 (29 Nov 2023)-img/p1_phy_s1_q02_1.png}{67.6}{0}
\qfigdef{p1_phy_s1_q03_1}{Pre-Score (2023-24) - Test 07 (29 Nov 2023)-img/p1_phy_s1_q03_1.png}{138.5}{0}
\qfigdef{p1_phy_s1_q04_1}{Pre-Score (2023-24) - Test 07 (29 Nov 2023)-img/p1_phy_s1_q04_1.png}{125.9}{0}
\qfigdef{p1_phy_s1_q05_1}{Pre-Score (2023-24) - Test 07 (29 Nov 2023)-img/p1_phy_s1_q05_1.png}{89.9}{0}
\qfigdef{p1_phy_s2_q07_1}{Pre-Score (2023-24) - Test 07 (29 Nov 2023)-img/p1_phy_s2_q07_1.png}{105.5}{0}
\qfigdef{p1_phy_s2_q08_1}{Pre-Score (2023-24) - Test 07 (29 Nov 2023)-img/p1_phy_s2_q08_1.png}{131.8}{0}
\qfigdef{p1_phy_s3_q01_1}{Pre-Score (2023-24) - Test 07 (29 Nov 2023)-img/p1_phy_s3_q01_1.png}{129.2}{0}
\qfigdef{p1_phy_s3_q02_1}{Pre-Score (2023-24) - Test 07 (29 Nov 2023)-img/p1_phy_s3_q02_1.png}{115.8}{0}
\qfigdef{p1_phy_s3_q03_1}{Pre-Score (2023-24) - Test 07 (29 Nov 2023)-img/p1_phy_s3_q03_1.png}{156.7}{0}
\qfigdef{p1_phy_s3_q04_1}{Pre-Score (2023-24) - Test 07 (29 Nov 2023)-img/p1_phy_s3_q04_1.png}{120.0}{0}
\qfigdef{p1_phy_s3_q05_1}{Pre-Score (2023-24) - Test 07 (29 Nov 2023)-img/p1_phy_s3_q05_1.png}{151.6}{0}
\qfigdef{p1_phy_s3_q08_1}{Pre-Score (2023-24) - Test 07 (29 Nov 2023)-img/p1_phy_s3_q08_1.png}{129.4}{0}
\qfigdef{p1_chem_s2_q09_1}{Pre-Score (2023-24) - Test 07 (29 Nov 2023)-img/p1_chem_s2_q09_1.png}{92.2}{0}
\qfigdef{p1_chem_s2_q09_2}{Pre-Score (2023-24) - Test 07 (29 Nov 2023)-img/p1_chem_s2_q09_2.png}{90.6}{0}
\qfigdef{p1_chem_s2_q09_3}{Pre-Score (2023-24) - Test 07 (29 Nov 2023)-img/p1_chem_s2_q09_3.png}{91.7}{0}
\qfigdef{p1_chem_s2_q09_4}{Pre-Score (2023-24) - Test 07 (29 Nov 2023)-img/p1_chem_s2_q09_4.png}{82.6}{0}
\qfigdef{p1_chem_s2_q09_5}{Pre-Score (2023-24) - Test 07 (29 Nov 2023)-img/p1_chem_s2_q09_5.png}{85.7}{0}
\qfigdef{p1_chem_s3_q10_2}{Pre-Score (2023-24) - Test 07 (29 Nov 2023)-img/p1_chem_s3_q10_2.png}{180.1}{0}
\qfigdef{p2_phy_s1_q02_1}{Pre-Score (2023-24) - Test 07 (29 Nov 2023)-img/p2_phy_s1_q02_1.png}{129.4}{0}
\qfigdef{p2_phy_s1_q04_1}{Pre-Score (2023-24) - Test 07 (29 Nov 2023)-img/p2_phy_s1_q04_1.png}{82.0}{0}
\qfigdef{p2_phy_s2_q05_1}{Pre-Score (2023-24) - Test 07 (29 Nov 2023)-img/p2_phy_s2_q05_1.png}{151.6}{0}
\qfigdef{p2_phy_s2_q06_1}{Pre-Score (2023-24) - Test 07 (29 Nov 2023)-img/p2_phy_s2_q06_1.png}{151.6}{0}
\qfigdef{p2_phy_s2_q07_1}{Pre-Score (2023-24) - Test 07 (29 Nov 2023)-img/p2_phy_s2_q07_1.png}{97.9}{0}
\qfigdef{p2_phy_s2_q08_1}{Pre-Score (2023-24) - Test 07 (29 Nov 2023)-img/p2_phy_s2_q08_1.png}{142.0}{0}
\qfigdef{p2_phy_s2_q09_1}{Pre-Score (2023-24) - Test 07 (29 Nov 2023)-img/p2_phy_s2_q09_1.png}{142.1}{0}
\qfigdef{p2_phy_s2_q10_1}{Pre-Score (2023-24) - Test 07 (29 Nov 2023)-img/p2_phy_s2_q10_1.png}{105.6}{0}
\qfigdef{p2_phy_s3_q01_1}{Pre-Score (2023-24) - Test 07 (29 Nov 2023)-img/p2_phy_s3_q01_1.png}{87.2}{0}
\qfigdef{p2_phy_s3_q02_1}{Pre-Score (2023-24) - Test 07 (29 Nov 2023)-img/p2_phy_s3_q02_1.png}{109.5}{0}
\qfigdef{p2_phy_s3_q03_1}{Pre-Score (2023-24) - Test 07 (29 Nov 2023)-img/p2_phy_s3_q03_1.png}{103.1}{0}
\qfigdef{p2_phy_s3_q05_1}{Pre-Score (2023-24) - Test 07 (29 Nov 2023)-img/p2_phy_s3_q05_1.png}{117.1}{0}
\qfigdef{p2_phy_s3_q07_1}{Pre-Score (2023-24) - Test 07 (29 Nov 2023)-img/p2_phy_s3_q07_1.png}{150.6}{0}
\qfigdef{p2_chem_s1_q03_1}{Pre-Score (2023-24) - Test 07 (29 Nov 2023)-img/p2_chem_s1_q03_1.png}{126.8}{0}
\qfigdef{p2_chem_s2_q07_1}{Pre-Score (2023-24) - Test 07 (29 Nov 2023)-img/p2_chem_s2_q07_1.png}{76.8}{0}
\qfigdef{p2_chem_s3_q01_2}{Pre-Score (2023-24) - Test 07 (29 Nov 2023)-img/p2_chem_s3_q01_2.png}{158.9}{0}
\raggedcolumns\makeatletter\@beginparpenalty=10000\makeatother

\begin{document}
\begin{center}{\Large\bfseries Pre-Score (2023-24) --- Test 07 (29 Nov 2023)}\end{center}
\begin{multicols}{2}
\papersec{1}
\partsec{1}{PHYSICS}
\sect{SECTION-I (i)}
\qn{1} For the given circuit, choose the CORRECT statements: (given all capacitors have capacitance $=1$ F \& all resistors have resistance $=10\ \Omega$)
\qfig{p1_phy_s1_q01_1}
\begin{qoptions}
\opt{A}{Potential difference between point A \& B ($V_A-V_B$) is 40 V.}
\opt{B}{Potential difference between point A \& B ($V_A-V_B$) is 50 V.}
\opt{C}{Magnitude of charge $Q_1$ is 30 C.}
\opt{D}{Magnitude of charge $Q_2$ is 50 C.}
\end{qoptions}
\qn{2} Two large square metallic plates ($\ell\times\ell$) are separated by a distance $d$. The potential of plates remain constant at $\dfrac V2$ \& $-\dfrac V2$. The plates are connected by a smooth horizontal insulating platform. A very small metallic ball of mass $m$ and radius $r$ can move on the platform. The collisions are perfectly inelastic. ($\ell\gg d\gg r$)
\qfig{p1_phy_s1_q02_1}
\begin{qoptions}
\opt{A}{The effective charge that flows from left plate to right plate in a complete round trip is $4\pi\varepsilon_0 Vr$.}
\opt{B}{The average current over a large period of time is proportional to $V^{3/2}$.}
\opt{C}{The ``resistance'' to current flow $\left(\dfrac Vi\right)$ is inversely proportional to $V$.}
\opt{D}{The average power loss is proportional to $V^2$.}
\end{qoptions}
\qn{3} A potentiometer is connected to the circuit as shown in figure. Length of wire AB is 200 cm and resistance is $8\,\Omega$.
\qfig{p1_phy_s1_q03_1}
\begin{qoptions}
\opt{A}{When only switch $S_1$ is closed null point is not obtained.}
\opt{B}{When only $S_2$ is closed null point is at 150 cm from A.}
\opt{C}{When both $S_1$ and $S_2$ are closed null point is not obtained.}
\opt{D}{When $S_1$, $S_2$ and $S_3$ all are closed null point is at 50 cm from A.}
\end{qoptions}
\qn{4} ABCDEFGH is a cube made up of wire frame, there is equal resistance $R$ in each side of cube. Four identical ideal ammeters are connected between vertices BE, EG, HC and FC as shown in figure. An ideal cell of EMF $V$ is connected between A \& F. Connections are only at vertices of cube. Choose correct readings of ammeters.
\qfig{p1_phy_s1_q04_1}
\begin{qoptions}
\opt{A}{Reading of $A_1$ is $3v/9R$}
\opt{B}{Reading of $A_2$ is $2v/9R$}
\opt{C}{Reading of $A_3$ is $7v/9R$}
\opt{D}{Reading of $A_4$ is $14v/9R$}
\end{qoptions}
\qn{5} Two concentric conducting thin shell of radii $a$ and $2a$ having charge $Q$ and $-Q$ respectively kept as shown in figure. Another point charge $Q$ also kept at a distance $\dfrac a2$ from common center, then
\qfig{p1_phy_s1_q05_1}
\begin{qoptions}
\opt{A}{Energy stored in electric field from $r=a$ to $r=2a$ is $\dfrac{Q^2}{4\pi\varepsilon_0 a}$, where $r$ is radial distance from common center.}
\opt{B}{Electric potential of inner shell is $\dfrac{3Q}{8\pi\varepsilon_0 a}$.}
\opt{C}{If switch S is closed, then charge flow through the switch into the earth is $\dfrac{5Q}{2}$.}
\opt{D}{Electric field at common center remains unchanged on closing the switch S.}
\end{qoptions}
\qn{6} A particle of mass $m$ and charge $(-q)$ is constrained to move in the horizontal plane. The particle can oscillate back and forth in a straight line. A charge particle ($Q\gg q$) is placed at a distance $R$ below the origin of the horizontal plane and another charge $(-q)$ is fixed at the origin of the horizontal plane.
\begin{qoptions}
\opt{A}{The distance from origin at which the particle is in equilibrium is $r_0=R\left(\dfrac qQ\right)^{1/3}$.}
\opt{B}{The distance from origin at which the particle is in equilibrium is $r_0=R\left(\dfrac{3q}{Q}\right)^{1/3}$.}
\opt{C}{The angular frequency of oscillations for small displacements around the equilibrium position is $\sqrt{\dfrac{qQ}{4\pi\epsilon_0 mR^3}}$.}
\opt{D}{The angular frequency of oscillations for small displacements around the equilibrium position is $\sqrt{\dfrac{3qQ}{4\pi\epsilon_0 mR^3}}$.}
\end{qoptions}
\sect{SECTION-I (ii)}
\qn{7} Initially the switch `S' is open for a long time, Now the switch `S' is closed at $t=0$ ($R_1=2\Omega$, $R_2=2\Omega$, $C_1=1$F and $C_2=1$F).
\qfig{p1_phy_s2_q07_1}
\par\nobreak
{\small\setlength{\tabcolsep}{3pt}\renewcommand{\arraystretch}{1.5}
\begin{tabular}{|L{0.8cm}|L{3.3cm}|L{0.8cm}|L{1.6cm}|}\hline
\multicolumn{2}{|c|}{\textbf{List-I}} & \multicolumn{2}{c|}{\textbf{List-II}}\\\hline
(I) & The charge on capacitor $C_1$ (in C) just before the switch is closed & (P) & $\dfrac{1}{2e}$\\\hline
(II) & The charge on capacitor $C_1$ (in C) just after the switch is closed & (Q) & $1-\dfrac{1}{2e}$\\\hline
(III) & The current through AB (in A) at $t=2$ sec is & (R) & 1\\\hline
(IV) & The charge on capacitor $C_1$ (in C) long time after the switch is closed is & (S) & $\dfrac12$\\\hline
 & & (T) & None\\\hline
\end{tabular}}
\begin{qoptions}
\opt{A}{I $\to$ R; II $\to$ R; III $\to$ Q; IV $\to$ R}
\opt{B}{I $\to$ R; II $\to$ S; III $\to$ P; IV $\to$ R}
\opt{C}{I $\to$ R; II $\to$ S; III $\to$ Q; IV $\to$ P}
\opt{D}{I $\to$ R; II $\to$ R; III $\to$ P; IV $\to$ S}
\end{qoptions}
\qn{8} All identical capacitor plates each of area $A$ are arranged such that adjacent plates are at $d$ distance apart. Plates are connected to a source of emf $V$ volts as shown in figure. Match the quantities in List-I with their respective answer in List-II. (dielectric constant in each shaded region is $k=2$)
\qfig{p1_phy_s2_q08_1}
\par\nobreak
{\small\setlength{\tabcolsep}{3pt}\renewcommand{\arraystretch}{1.6}
\begin{tabular}{|L{0.8cm}|L{3.5cm}|L{0.8cm}|L{1.0cm}|}\hline
\multicolumn{2}{|c|}{\textbf{List-I}} & \multicolumn{2}{c|}{\textbf{List-II}}\\\hline
(I) & Charge on plate 3 (multiple of $\dfrac{\varepsilon_0AV}{d}$) & (P) & 4/3\\\hline
(II) & Energy stored in the system (multiple of $\dfrac{13\varepsilon_0AV^2}{d}$) & (Q) & 1/3\\\hline
(III) & Equivalent capacitance between AB (multiple of $\dfrac{13A\varepsilon_0}{d}$) & (R) & 2/3\\\hline
(IV) & Potential difference between the plates 4 and 7 (multiple of $V/2$) & (S) & 5/3\\\hline
 & & (T) & None\\\hline
\end{tabular}}
\begin{qoptions}
\opt{A}{I $\to$ S; II $\to$ R; III $\to$ P; IV $\to$ Q}
\opt{B}{I $\to$ P; II $\to$ Q; III $\to$ S; IV $\to$ R}
\opt{C}{I $\to$ R; II $\to$ S; III $\to$ P; IV $\to$ Q}
\opt{D}{I $\to$ S; II $\to$ Q; III $\to$ R; IV $\to$ P}
\end{qoptions}
\qn{9} Answer the following by appropriately matching the lists based on the information given in the paragraph. A particle is projected vertically up from the surface of the earth with a speed $\sqrt{gR}$. Which of the following is the only correct combination?
\par\nobreak
{\small\setlength{\tabcolsep}{3pt}\renewcommand{\arraystretch}{1.5}
\begin{tabular}{|L{0.8cm}|L{3.2cm}|L{0.8cm}|L{2.1cm}|}\hline
\multicolumn{2}{|c|}{\textbf{List-I}} & \multicolumn{2}{c|}{\textbf{List-II}}\\\hline
(I) & Maximum height attained by the particle is $\alpha R$ the value of $\alpha$ is & (P) & 1\\\hline
(II) & The duration for which the particle is between the point of projection and half of maximum height (including upward and downward journey) is $\beta\sqrt{\dfrac{R^3}{GM}}$, the value of $\beta$ is & (Q) & 2\\\hline
(III) & Time taken by the particle to travel from half of maximum height to maximum height is $\gamma\sqrt{\dfrac{R^3}{GM}}$, the value of $\gamma$ is & (R) & $\dfrac{\pi}{3}+\dfrac{\sqrt3}{2}$\\\hline
(IV) & The acceleration due to gravity at maximum height is $\dfrac{\delta GM}{4R^2}$, the value of $\delta$ is & (S) & $\left(\dfrac{\pi}{3}-\sqrt3+2\right)$\\\hline
 & & (T) & None\\\hline
\end{tabular}}
\begin{qoptions}
\opt{A}{I $\to$ P; II $\to$ S; III $\to$ R; IV $\to$ P}
\opt{B}{I $\to$ P; II $\to$ Q; III $\to$ S; IV $\to$ R}
\opt{C}{I $\to$ R; II $\to$ P; III $\to$ R; IV $\to$ P}
\opt{D}{I $\to$ P; II $\to$ R; III $\to$ S; IV $\to$ P}
\end{qoptions}
\qn{10} Electric field in a region is given by $\vec E=\left[(2xy)\hat i+(x^2)\hat j\right]$ V/m. Electric potential at origin is 5V. Match the following two columns.
\par\nobreak
{\small\setlength{\tabcolsep}{3pt}\renewcommand{\arraystretch}{1.3}
\begin{tabular}{|L{0.8cm}|L{3.3cm}|L{0.8cm}|L{1.0cm}|}\hline
\multicolumn{2}{|c|}{\textbf{List-I}} & \multicolumn{2}{c|}{\textbf{List-II}}\\\hline
(I) & Potential at point (5m, 2m, 1m) & (P) & 4V\\\hline
(II) & Potential at point (1m, 1m, 1m) & (Q) & 1V\\\hline
(III) & Potential at point (2m, 1m, 1m) & (R) & 3V\\\hline
(IV) & Potential at point (1m, 2m, 1m) & (S) & $-45$V\\\hline
 & & (T) & None\\\hline
\end{tabular}}
\begin{qoptions}
\opt{A}{I $\to$ S; II $\to$ P; III $\to$ R; IV $\to$ Q}
\opt{B}{I $\to$ R; II $\to$ P; III $\to$ S; IV $\to$ Q}
\opt{C}{I $\to$ R; II $\to$ S; III $\to$ P; IV $\to$ Q}
\opt{D}{I $\to$ S; II $\to$ P; III $\to$ Q; IV $\to$ R}
\end{qoptions}
\sect{SECTION-II}
\qn{1} The circuit shows the pattern of ``infinite'' connection of resistances. Each resistance has a value $R=1\Omega$. The equivalent resistance between A and B is
\qfig{p1_phy_s3_q01_1}
\qn{2} A spherical cavity of radius $\dfrac R4$ is cut from a sphere of radius $R$ of density $\rho_1$. The center of cavity is at a distance $\dfrac R2$ from the center of the sphere. The cavity is filled with a material of density $\rho_2$, as shown in diagram. The acceleration due to gravity at point A is 10 \% more than acceleration due to gravity at point B, the value of $\dfrac{\rho_1}{\rho_2}$ is
\qfig{p1_phy_s3_q02_1}
\qn{3} Separation between the plates of capacitor is $d$ and cross section area of plates is $Lb$. Dielectric constant is 4. The dielectric is connected to a block of mass $m$ as shown in the diagram. The value of $m$ (in kg) for which the block moves with a constant velocity is. Given : $\dfrac{\varepsilon_0 b}{d}=10$ (in SI unit)
\qfig{p1_phy_s3_q03_1}
\qn{4} A cone with base radius $\sqrt3R$ \& height $R$ has charge $q$ and $2q$ placed at vertices \& base centre respectively. The total electric flux passing through cone is $x\dfrac{q}{\varepsilon_0}$. Find $x$.
\qfig{p1_phy_s3_q04_1}
\qn{5} A long cylindrical capacitor has the cross section shown in figure. The cylinder of radius $R$ has charge density $\sigma$. Rest are uncharged. Find $V_{AB}/V_{BC}$.
\qfig{p1_phy_s3_q05_1}
\qn{6} The range of ammeter can be changed by connecting a resistance box ABCD across the galvanometer. The terminals AD are connected across the galvanometer terminals. When the terminals AC are shorted, the range of ammeter is 4A. When terminals BD are shorted, the range of ammeter becomes 3A. When terminals BC are shorted, the range of ammeter becomes 2A. The full deflection current of the galvanometer is
\qn{7} Two small dipoles having dipole moments $p\,\hat i$ and $\sqrt3 p\,\hat j$ are placed at the origin. The point on the circle of radius $R$, where the potential is maximum is $(x,y)$. The ratio $\dfrac yx$ is
\qn{8} Four resistance $R$, $3R$, $R$ and $3R$ are connected as shown in figure. The two voltmeters are identical and Ammeter is ideal. The readings of ammeter and voltmeter are shown in the diagram. The value of $R$ is.
\qfig{p1_phy_s3_q08_1}
\partsec{2}{CHEMISTRY}
\sect{SECTION-I (i)}
\qn{1} Which of the following statement is/are CORRECT?
\begin{qoptions}
\opt{A}{Higher is the gold number of a lyophilic sol more is its coagulating power}
\opt{B}{An electrolyte can coagulate both positive sol and negative sol}
\opt{C}{All impurities present in a colloid makes it more stable}
\opt{D}{Dialysis is a process to remove impurities of ions and smaller molecules from a colloid}
\end{qoptions}
\qn{2} Kolbe's electrolysis of sodium acetate was carried out in aqueous medium. Given 75\% yield of electrolysis process 4.4 gm CO$_2$ is produced at anode. Select correct option(s) :
\begin{qoptions}
\opt{A}{Total charge passed in 8/3F}
\opt{B}{Mass of hydrocarbon produced is 1.5 gm}
\opt{C}{Total mass of gases produced during the process is 6.1 gm}
\opt{D}{Total mass of gases produced during the process is 6 gm}
\end{qoptions}
\qn{3} At 298 K, the standard reduction potentials are 1.51 V for MnO$_4^{\ominus}|$Mn$^{+2}$, 1.36 V for Cl$_2|$Cl$^{\ominus}$, 1.07 V for Br$_2\,|\,$Br$^{\ominus}$, and 0.54 V for I$_2|$I$^{\ominus}$. At pH = 3, permanganate is expected to oxidize $\left(\dfrac{RT}{F}=0.059\ \text{V}\right)$ :$-$
\begin{qoptions}
\opt{A}{Cl$^{\ominus}$}
\opt{B}{Br$^{\ominus}$}
\opt{C}{I$^{\ominus}$}
\opt{D}{All Cl$^{\ominus}$, Br$^{\ominus}$, I$^{\ominus}$}
\end{qoptions}
\qn{4} Which is/are correct statement(s)
\begin{qoptions}
\opt{A}{A catalyst may appear in the kinetic rate law equation of the reaction}
\opt{B}{A reaction having first order may be either elementary or complex.}
\opt{C}{A plot of rate of reaction v/s concentration of reactant for first order reaction (A $\to$ Product) is a straight line with positive slope}
\opt{D}{A zero order reaction may be elementary or complex.}
\end{qoptions}
\qn{5} Which of the following statement(s) is(are) correct ?
\begin{qoptions}
\opt{A}{Smoke and cloud both are aerosols with same dispersion medium.}
\opt{B}{Surface tension of lyophilic sols is less than the surface tension of dispersion medium while that of lyophobic sols is more than that of dispersion medium.}
\opt{C}{Lyophobic sols can be observed using ultra microscope}
\opt{D}{Colloidal sol shows low magnitude of colligative properties, as compared to true solutions.}
\end{qoptions}
\qn{6} For the reaction at 273\textdegree C\\
N$_2$(g) + 3H$_2$(g) $\to$ 2NH$_3$(g)\\
It is observed that the pressure of H$_2$(g) falls from 760 mmHg to 380 mmHg in 100 sec.\\
Select the correct statement(s).
\begin{qoptions}
\opt{A}{Average rate of reaction is $5\times10^{-3}$ atm/sec.}
\opt{B}{Average rate of formation of NH$_3$(g) is $\dfrac{1}{300}$ atm/sec.}
\opt{C}{Average rate of disappearance of N$_2$(g) is twice of rate of appearance of NH$_3$(g).}
\opt{D}{$\log\left(\dfrac{-\Delta P_{H_2}}{\Delta t}\right)=\log\left(\dfrac{-\Delta P_{N_2}}{\Delta t}\right)+0.4771$}
\end{qoptions}
\end{multicols}
\noindent\begin{minipage}{\textwidth}
\sect{SECTION-I (ii)}
\qn{7}\par\nobreak
\small\setlength{\tabcolsep}{4pt}\renewcommand{\arraystretch}{1.5}
\begin{center}\begin{tabular}{|L{0.8cm}|L{6.5cm}|L{0.8cm}|L{5.0cm}|}\hline
\multicolumn{2}{|c|}{\begin{tabular}[c]{@{}c@{}}\textbf{List-I}\\\textbf{Electrode (25\textdegree\,C 1atm)}\end{tabular}} & \multicolumn{2}{c|}{\begin{tabular}[c]{@{}c@{}}\textbf{List-II}\\\textbf{Oxidation potential (approximate)}\end{tabular}}\\\hline
(P) & Pt, H$_2|$NaCl(aq) & (1) & 0.66 volts\\\hline
(Q) & Pt, H$_2|$ NH$_4$Cl (1M) ($K_b$ of NH$_3=10^{-6}$) & (2) & 0.6 volts\\\hline
(R) & Pt, H$_2|$ HCOOK (1M) ($K_b$ of HCOO$^-=10^{-6}$) & (3) & 0.24 volts\\\hline
(S) & Pt, H$_2|$ KSH ($K_{a_1}=10^{-4}$, $K_{a_2}=10^{-10}$) & (4) & 0.42 volts\\\hline
 & & (5) & 0.36 volt\\\hline
\end{tabular}\end{center}
[Given : $\dfrac{2.303\,RT}{F}=0.06$]\\
The correct match is :
\begin{qoptions}
\opt{A}{P $\to$ 4; Q $\to$ 4; R $\to$ 1; S $\to$ 5}
\opt{B}{P $\to$ 4; Q $\to$ 2; R $\to$ 3; S $\to$ 4}
\opt{C}{P $\to$ 4; Q $\to$ 3; R $\to$ 2; S $\to$ 2}
\opt{D}{P $\to$ 4; Q $\to$ 3; R $\to$ 1; S $\to$ 4}
\end{qoptions}
\end{minipage}

\medskip
\noindent\begin{minipage}{\textwidth}
\qn{8} Match the List-I with List-II.\par\nobreak
\small\setlength{\tabcolsep}{4pt}\renewcommand{\arraystretch}{1.4}
\begin{center}\begin{tabular}{|L{0.8cm}|L{6.2cm}|L{0.8cm}|L{6.2cm}|}\hline
\multicolumn{2}{|c|}{\textbf{List-I}} & \multicolumn{2}{c|}{\textbf{List-II}}\\\hline
(P) & Molten PbCl$_2$ using inert electrode & (1) & Metal of salt will be reduced\\\hline
(Q) & NaCl aqueous solution using inert electrode & (2) & 2H$_2$O + 2e$^-$ $\to$ H$_2$(g) + 2OH$^-$\\\hline
(R) & Silver nitrate aqueous solution using Ag electrodes & (3) & Solution becomes basic after electrolysis\\\hline
(S) & Sodium nitrate aqueous solution using inert electrodes & (4) & Solution becomes acidic after electrolysis\\\hline
 & & (5) & Solution becomes neutral after electrolysis\\\hline
\end{tabular}\end{center}
[Ignore hydrolysis of ions in solution]
\begin{qoptions}
\opt{A}{P $\to$ 1; Q $\to$ 3; R $\to$ 5; S $\to$ 2}
\opt{B}{P $\to$ 1; Q $\to$ 2; R $\to$ 1; S $\to$ 3}
\opt{C}{P $\to$ 5; Q $\to$ 5; R $\to$ 4; S $\to$ 5}
\opt{D}{P $\to$ 5; Q $\to$ 2; R $\to$ 4; S $\to$ 5}
\end{qoptions}
\end{minipage}

\medskip
\noindent\begin{minipage}{\textwidth}
\qn{9} Match the rate expressions in LIST-I for the decomposition of X with the corresponding profiles provided in LIST-II. $X_s$ and $k$ constants having appropriate units.\par\nobreak
\small\setlength{\tabcolsep}{4pt}\renewcommand{\arraystretch}{1.3}
\begin{center}\begin{tabular}{|L{0.8cm}|L{6.0cm}|L{0.8cm}|L{4.5cm}|}\hline
\multicolumn{2}{|c|}{\textbf{LIST-I}} & \multicolumn{2}{c|}{\textbf{LIST-II}}\\\hline
(P) & rate $=\dfrac{k[X]}{X_s+[X]}$\newline under all possible initial concentration of X & (1) & \qfigi{p1_chem_s2_q09_1}\\\hline
(Q) & rate $=\dfrac{k[X]}{X_s+[X]}$\newline where initial concentration of X are much less than $X_s$ & (2) & \qfigi{p1_chem_s2_q09_2}\\\hline
(R) & rate $=\dfrac{k[X]}{X_s+[X]}$\newline where initial concentration of X are much higher than $X_s$ & (3) & \qfigi{p1_chem_s2_q09_3}\\\hline
(S) & rate $=\dfrac{k[X]^2}{X_s+[X]}$\newline where initial concentration of X is much higher than $X_s$ & (4) & \qfigi{p1_chem_s2_q09_4}\\\hline
 & & (5) & \qfigi{p1_chem_s2_q09_5}\\\hline
\end{tabular}\end{center}
\begin{qoptions}
\opt{A}{P $\to$ 3; Q $\to$ 2; R $\to$ 4; S $\to$ 5}
\opt{B}{P $\to$ 3; Q $\to$ 4; R $\to$ 4; S $\to$ 5}
\opt{C}{P $\to$ 1; Q $\to$ 2; R $\to$ 2; S $\to$ 3}
\opt{D}{P $\to$ 1; Q $\to$ 4; R $\to$ 2; S $\to$ 3}
\end{qoptions}
\end{minipage}
\begin{multicols}{2}
\qn{10} Which of the following option show correct matches :-
\par\nobreak
{\small\setlength{\tabcolsep}{3pt}\renewcommand{\arraystretch}{1.3}
\begin{tabular}{|L{0.8cm}|L{3.0cm}|L{0.8cm}|L{2.7cm}|}\hline
\multicolumn{2}{|c|}{\textbf{List-I}} & \multicolumn{2}{c|}{\textbf{List-II}}\\\hline
(P) & Precipitation of colloidal sol & (1) & Peptization\\\hline
(Q) & Purification of colloid & (2) & Coagulation\\\hline
(R) & Gold Number & (3) & Desorption\\\hline
(S) & Formation of colloidal sol & (4) & Dialysis\\\hline
 & & (5) & Measure of protective power\\\hline
\end{tabular}}
\begin{qoptions}
\opt{A}{P $\to$ 1; Q $\to$ 4; R $\to$ 3; S $\to$ 1}
\opt{B}{P $\to$ 5; Q $\to$ 2; R $\to$ 1; S $\to$ 3}
\opt{C}{P $\to$ 2; Q $\to$ 4; R $\to$ 3; S $\to$ 1}
\opt{D}{P $\to$ 2; Q $\to$ 4; R $\to$ 5; S $\to$ 1}
\end{qoptions}
\sect{SECTION-II}
\qn{1} Consider the following statements:\\
1. Osmotic pressure of rubber sol will be same as that of sucrose solution having same mass of rubber and sucrose mixed in same mass of H$_2$O\\
2. Soap sol of sodium palmitate will be coagulated near cathode on electrophoresis.\\
3. For miscellisation of surfactants entropy of system increases.\\
4. Movement of dispersion medium under the influence of an electric field is called cataphoresis.\\
5. Gold sol is a negatively charged hydrophilic sol.\\
6. Tyndall effect is more distinct for lyophobic colloid than for lyophilic colloids.\\
7. Gold sol on mixing with starch sol causes stablisation of starch sol.\\
8. During adsorption there will be an increase in residual forces on the molecules.\\
9. The extent of adsorption per unit mass of adsorbent will remain same even on changing surface area of same mass of adsorbent.\\
10. On adding 0.1 M KI solution slowly to 0.01 M AgNO$_3$ solution, causes conductivity to increase gradually and after some time it becomes constant.\\
If number of correct statements $=x$\\
And number of incorrect statements $=y$\\
Then, $\dfrac xy=$
\qn{2} If ``X'' represents reduction potential ($E^0$, in V) of MnO$_4^-$(aq)/Mn(s) and ``Y'' represents reduction potential ($E^0$, in V) of Cr$_2$O$_7^{2-}$(aq.)/Cr(s) then $|X-Y|$ is.\\
{\raggedright {[}Given : $E^0_{(\mathrm{MnO_4^-(aq)/MnO_2(s)})}=1.7$V ; $E^0_{(\mathrm{MnO_2(s)/Mn^{2+}(aq)})}=1.25$V ; $E^0_{(\mathrm{Mn^{2+}(aq)/Mn(s)})}=-1$V ]\\
{[}$E^0_{(\mathrm{Cr_2O_7^{2-}(aq)/Cr^{3+}(aq)})}=1.36$V ; $E^0_{(\mathrm{Cr^{3+}(aq)/Cr(s)})}=-0.74$V ]\par}
\qn{3} The conductivity of 0.001M Na$_2$SO$_4$ solution is $2.6\times10^{-4}$ S cm$^{-1}$ and it increases to $8.2\times10^{-4}$ S cm$^{-1}$, when the solution is saturated with CaSO$_4$. The limiting molar conductances of Na$^+$ and Ca$^{2+}$ ions are 50 and 120 S cm$^2$ mol$^{-1}$, respectively. Neglecting conductivity of water, the solubility product of CaSO$_4$ is $Y\times10^{-x}$ M$^2$. Then $Y$ is.
\qn{4} The decomposition of A$_2$B$_2$ follows the rate equation $r=K[\mathrm{A_2B_2}]^2$.\\
A$_2$B$_2$ $\xrightarrow{K}$ A$_2$ + B$_2$. The kinetic data of the reaction is given below.\par\nobreak
{\small\renewcommand{\arraystretch}{1.3}\begin{tabular}{|L{4.0cm}|L{1.6cm}|}\hline
Initial conc. of A$_2$B$_2$ (mole/lit) & 7\\\hline
First half life of A$_2$B$_2$ in sec. & 1440\\\hline
\end{tabular}}\par
If the initial concentration of A$_2$B$_2$ were 56 m mole/lit. and the time taken for the conversion of A$_2$B$_2$ to reach 60\% of initial value is `t' seconds then find value of $\left(\dfrac{t}{1000}\right)$?
\qn{5} If mass of copper deposited in a copper voltameter is 30 grams in 30 seconds. The time current graph is as shown in figure, calculate electrochemical equivalent (gm / coulomb) of Copper during the entire process.
\qfig{p1_chem_s3_q10_2}
\qn{6} For any acid catalysed reaction\\
A $\xrightarrow{H^+}$ B\\
half life period is independent of of concentration of A at given pH. At given concentration of A, half life period is 10 min at pH $=2$ and half life time is 1000 min at pH $=3$. If the rate law expression of reaction is :\\
Rate $=K[A]^x[H^+]^y$, then calculate value of $(x+y)$?
\qn{7} For the first order reaction
\[\mathrm{A(aq)\to 2B(aq)+C(aq)};\]
where A is an optically active compound while B is reducing agent which reacts with H$_2$O$_2$. The progress of reaction is monitored by measuring optical rotation as well by titrating a definite volume of reaction mixture with H$_2$O$_2$ from time to time. In an experiment, optical rotation was found to be $\theta=30^\circ$ at $t=20$ min and $\theta=15^\circ$ at time $t=50$ min from the start of reaction. If in the same experiment volume of H$_2$O$_2$ consumed at $t=30$ min from start is 30 mL then volume of H$_2$O$_2$ which is consumed by same volume of reaction mixture at $t=90$ min will be :\\
{[}Compound `B' and `C' are optically inactive]
\qn{8} 10\% sites of catalyst bed has been occupied by H$_2$ gas molecules. On heating H$_2$ gas is evolved from sites and collected at 0.030 atm and 300K in a vessel of volume 2.46 cm$^3$. Total number of sites available on catalyst is $6\times10^{15}$ per cm$^2$ and surface area of catalyst is $10^3$cm$^2$. The number of H$_2$ molecules adsorbed per site of catalyst are :\\
(Given $N_A=6\times10^{23}$, $R=0.082$ L.atm K$^{-1}$ mol$^{-1}$)
\partsec{3}{MATHEMATICS}
\sect{SECTION-I (i)}
\qn{1} $f:[0,1]\to\mathbb R$ be a function such that $(f(x))^2=1$ and $f(x)$ is discontinuous at exactly two points and $\displaystyle\int_0^1 f(x)\,dx=0$. Let $I=\displaystyle\int_0^1 xf(x)\,dx$, choose the incorrect option(s)
\begin{qoptions}
\opt{A}{maximum value of $I$ is $\dfrac14$}
\opt{B}{minimum value of $I$ is $-\dfrac34$}
\opt{C}{maximum value of $I$ is 3}
\opt{D}{minimum value of $I$ is $-2$}
\end{qoptions}
\qn{2} Let $S_n=\left(\dfrac1{n^6}+\dfrac{32}{n^6}+\dfrac{243}{n^6}+\dfrac{1024}{n^6}+\ldots+\dfrac1n\right)$ and $T_n=\left(\dfrac1{n^6}+\dfrac{32}{n^6}+\dfrac{243}{n^6}+\dfrac{1024}{n^6}+\ldots+\dfrac{(n-1)^5}{n^6}\right)$, then which of the following is/are true ?
\begin{qoptions}
\opt{A}{$S_n>\dfrac16$}
\opt{B}{$(S_n+T_n)<\dfrac13$}
\opt{C}{$(S_n+T_n)>\dfrac13$}
\opt{D}{$T_n<\dfrac16$}
\end{qoptions}
\qn{3} The solution of $\dfrac{dy}{dx}=\dfrac{x^2+y^2+1}{2xy}$ satisfying $y(1)=1$ is given by :
\begin{qoptions}
\opt{A}{a hyperbola}
\opt{B}{a circle}
\opt{C}{$y^2=x(1+x)-1$}
\opt{D}{$(x-2)^2+(y-3)^2=5$}
\end{qoptions}
\qn{4} If $\displaystyle\int_2^3\frac{e^{x-1}}{(x^2-5x+4)}\,2x\,dx=AF(1)+BF(-2)$ and $F(t)=\displaystyle\int_t^{t+1}\frac{e^x}{x}\,dx$, then
\begin{qoptions}
\opt{A}{$A=-2/3$}
\opt{B}{$B=(4/3)e^3$}
\opt{C}{$A=2/3$}
\opt{D}{$B=(8/3)e^3$}
\end{qoptions}
\qn{5} Let $R$ be the region in the first quadrant bounded by the $x$-axis, the line $2y=x$, and the ellipse $\dfrac{x^2}{9}+y^2=1$. Let $R'$ be the region in the first quadrant bounded by the $y$-axis, the line $y=mx$, and the given ellipse, then the value of $m$ such that $R$ and $R'$ have the same area is
\begin{qoptions}
\opt{A}{$\dfrac29$}
\opt{B}{$\dfrac3{11}$}
\opt{C}{$\dfrac25$}
\opt{D}{$\dfrac27$}
\end{qoptions}
\qn{6} If a function $y=f(x)$ is such that $f'(x)$ is continuous function and satisfies $(f(x))^2=K+\displaystyle\int_0^x\left((f(t))^2+(f'(t))^2\right)dt$, $K\in\mathbb R^+$, then which of the following is/are possible ?
\begin{qoptions}
\opt{A}{$f(x)$ is increasing function $\forall\ x\in\mathbb R$}
\opt{B}{$f(x)$ is bounded function}
\opt{C}{$f(x)$ is neither even nor odd}
\opt{D}{if $K=100$ then $f(0)=10$}
\end{qoptions}
\fwstart
\sect{SECTION-I (ii)}
\qn{7} A function $F$ is defined by $F(x)=\displaystyle\int_1^x\frac{e^t}{t}\,dt\ \forall\ x>0$. Now express the functions in List-I in terms of $F$, then match the following List-I with List-II\par\nobreak
\begin{center}\begin{tabular}{|L{0.8cm}|L{4.5cm}|L{0.8cm}|L{5.5cm}|}\hline
\multicolumn{2}{|c|}{\textbf{List-I}} & \multicolumn{2}{c|}{\textbf{List-II}}\\\hline
(I) & $\displaystyle\int_1^x\frac{e^t}{t+2}\,dt$ & (P) & $F(x)-\dfrac{e^x}{x}+e$\\\hline
(II) & $\displaystyle\int_1^x\frac{e^{3t}}{t}\,dt$ & (Q) & $xe^{1/x}-e-F\!\left(\dfrac1x\right)$\\\hline
(III) & $\displaystyle\int_1^x\frac{e^t}{t^2}\,dt$ & (R) & $e^{-2}[F(x+2)-F(3)]$\\\hline
(IV) & $\displaystyle\int_1^x e^{1/t}\,dt$ & (S) & $F(3x)-F(3)$\\\hline
 & & (T) & $F(3x)+F(3)$\\\hline
\end{tabular}\end{center}
\begin{qoptions}
\opt{A}{I $\to$ S; II $\to$ R; III $\to$ Q; IV $\to$ P}
\opt{B}{I $\to$ R; II $\to$ S; III $\to$ P; IV $\to$ Q}
\opt{C}{I $\to$ R; II $\to$ S; III $\to$ Q; IV $\to$ P}
\opt{D}{I $\to$ P; II $\to$ Q; III $\to$ R; IV $\to$ S}
\end{qoptions}
\fwend
\fwstart
\qn{8} Match the following List-I with List-II (C is integration constant)\par\nobreak
\begin{center}\begin{tabular}{|L{0.8cm}|L{5.0cm}|L{0.8cm}|L{6.0cm}|}\hline
\multicolumn{2}{|c|}{\textbf{List-I}} & \multicolumn{2}{c|}{\textbf{List-II}}\\\hline
(I) & $\displaystyle\int\frac{x^4-1}{x^2\sqrt{x^4+x^2+1}}\,dx$ & (P) & $\ell n\left(\dfrac{(x^2+1)+\sqrt{x^4+1}}{x}\right)+C$\\\hline
(II) & $\displaystyle\int\frac{x^2-1}{x\sqrt{1+x^4}}\,dx$ & (Q) & $C+\dfrac1{\sqrt2}\ell n\left(\dfrac{\sqrt{x^4+1}-\sqrt2x}{(x^2-1)}\right)$\\\hline
(III) & $\displaystyle\int\frac{1+x^2}{(x^2-1)\sqrt{1+x^4}}\,dx$ & (R) & $C-\tan^{-1}\left(\sqrt{\sqrt{1+\dfrac1{x^4}}-1}\right)$\\\hline
(IV) & $\displaystyle\int\frac{dx}{(1+x^4)\sqrt{\sqrt{1+x^4}-x^2}}$ & (S) & $\dfrac{\sqrt{x^4+x^2+1}}{x}+C$\\\hline
 & & (T) & $\ell n\left(\dfrac{(x^2+1)-\sqrt{x^4+1}}{x}\right)+C$\\\hline
\end{tabular}\end{center}
\begin{qoptions}
\opt{A}{I $\to$ P; II $\to$ Q; III $\to$ R; IV $\to$ S}
\opt{B}{I $\to$ S; II $\to$ P; III $\to$ Q; IV $\to$ R}
\opt{C}{I $\to$ R; II $\to$ Q; III $\to$ P; IV $\to$ S}
\opt{D}{I $\to$ P; II $\to$ R; III $\to$ S; IV $\to$ Q}
\end{qoptions}
\fwend
\fwstart
\qn{9} Match the following List-I with List-II\par\nobreak
\begin{center}\begin{tabular}{|L{0.8cm}|L{8.0cm}|L{0.8cm}|L{3.0cm}|}\hline
\multicolumn{2}{|c|}{\textbf{List-I}} & \multicolumn{2}{c|}{\textbf{List-II}}\\\hline
(I) & $x\,dy=y(dx+y\,dy)$, $y(1)=1$ and $y(x_0)=-3$, then $x_0=\lambda$, then $\lambda$ is less than & (P) & $\dfrac14$\\\hline
(II) & If $y(t)$ is solution of $(t+1)\dfrac{dy}{dt}-ty=1$, $y(0)=-1$, and $y(1)=\lambda$, then $\lambda$ is greater than & (Q) & $-15$\\\hline
(III) & $(x^2+y^2)dy=xy\,dx$, $y(1)=1$, $y(x_0)=\sqrt e$ and $x_0=\lambda$, then $\lambda$ is greater than & (R) & $-\dfrac12$\\\hline
(IV) & $\dfrac{dy}{dx}+\dfrac{2y}{x}=0$, $y(1)=1$, then $y(2)$ is equal to & (S) & $\sqrt3\,e$\\\hline
 & & (T) & 1\\\hline
\end{tabular}\end{center}
\begin{qoptions}
\opt{A}{I $\to$ Q; II $\to$ R; III $\to$ S; IV $\to$ Q}
\opt{B}{I $\to$ Q; II $\to$ S; III $\to$ R; IV $\to$ P}
\opt{C}{I $\to$ Q; II $\to$ R; III $\to$ S; IV $\to$ P}
\opt{D}{I $\to$ P; II $\to$ S; III $\to$ R; IV $\to$ Q}
\end{qoptions}
\fwend
\fwstart
\qn{10} Match List-I with List-II :\par\nobreak
\begin{center}\begin{tabular}{|L{0.8cm}|L{6.3cm}|L{0.8cm}|L{4.2cm}|}\hline
\multicolumn{2}{|c|}{\textbf{List-I}} & \multicolumn{2}{c|}{\textbf{List-II}}\\\hline
(I) & The area bounded by $2^{|x|}|y|+2^{|x|-1}\le1$ is & (P) & $4-\dfrac{2}{\ell n2}$\\\hline
(II) & The area bounded by $2^{|x|}|y|+2^{|x|-1}\le1$, $|x|\le\dfrac12$ \& $|y|\le\dfrac12$ is equal to & (Q) & $\dfrac{2}{\ell n2}-2$\\\hline
(III) & The area bounded by $2^{|x|}|y|+2^{2|x|-1}\ge1$ \& $|x|+|y|\le1$ is equal to & (R) & $\dfrac{4\left(2^{1/3}-1\right)}{2^{1/3}\ell n2}-\dfrac23$\\\hline
(IV) & The area bounded by $2^{|x|}|y|+2^{|x|-1}\le1$ \& $|x|\le\dfrac13$ & (S) & $\dfrac{\left(\sqrt2-1\right)2\sqrt2}{\ell n2}-1$\\\hline
 & & (T) & $\dfrac{2}{\ell n2}+2$\\\hline
\end{tabular}\end{center}
\begin{qoptions}
\opt{A}{I $\to$ P; II $\to$ Q; III $\to$ R; IV $\to$ S}
\opt{B}{I $\to$ R; II $\to$ S; III $\to$ Q; IV $\to$ P}
\opt{C}{I $\to$ Q; II $\to$ S; III $\to$ P; IV $\to$ R}
\opt{D}{I $\to$ S; II $\to$ P; III $\to$ R; IV $\to$ Q}
\end{qoptions}
\fwend
\sect{SECTION-II}
\qn{1} If $f(x+f(y))+f(x-f(y))=f(x)\ \forall\ x,y\in\mathbb R$ and $(0,2)$ lies on $y=f(x)$, then $\displaystyle\int_0^{12}f(x)\,dx=\lambda\int_{-24}^{24}f(x-36)\,dx$ where $\lambda$ equals
\qn{2} If $\displaystyle\sum_{r=0}^{\infty}\left(\dfrac1{4r+1}-\dfrac1{4r+3}\right)=L$, then $[4L]$ is (where $[\,.\,]$ denotes greatest integer function)
\qn{3} If $f$ and $g$ are differentiable functions such that $xg(f(x))f'(g(x))g'(x)=f(g(x))g'(f(x))f'(x)$ for all real $x$. Also $f$ is non-negative and $g$ is positive. If $\displaystyle\int_0^p f(g(x))\,dx=\dfrac12-\dfrac{e^{-2p}}{2}$ for all real $p$ and $g(f(0))=1$. If the value of $g(f(4))$ is $e^{-k}$, where $k\in N$. Then $k$ is
\qn{4} If the line $y=mx$ bisects the area enclosed by the lines $x=0$, $y=0$, $x=\dfrac32$ and the curve $y=1+4x-x^2$ for $m=\dfrac pq$ ($p,q\in N$, $p$ \& $q$ are co-prime), then the value of $2p-3q=$
\qn{5} If $\displaystyle\int\frac{(x+2)\,dx}{\sqrt{x^4+8x^3+20x^2+16x+4}}=\dfrac12\ell n f(x)+c$, then $f(1)$ is equal to
\qn{6} Area of points satisfying $|2x+3y|+|3x-4y|\le2$ equals
\qn{7} Given ellipse $E:\dfrac{x^2}{a^2}+\dfrac{y^2}{b^2}=1$ and circle $C:(x-c)^2+y^2=c^2$. If $E$ is of smallest possible area ($\Delta$) such that it contains $C$, and $c=\dfrac{\sqrt2}{\sqrt{\sqrt3\pi}}$, then find $\Delta$.
\qn{8} Let $f$ be a continuous function that $p(x)=\displaystyle\int_{-2}^2 f(t)\,|x-t|\,dt$ where $x\in(-2,2)$, then the value of $\displaystyle\int_0^2\frac{p''(x)}{f(x)}\,dx$ is equal to
\end{multicols}
\papersec{2}
\begin{multicols}{2}
\partsec{1}{PHYSICS}
\sect{SECTION-I (i)}
\qn{1} A leaky capacitor is one which does not have a perfectly insulating dielectric layer between the plates upon which the charge is stored. Such a capacitor contains a dielectric material filling the space between the plates, and with an effective resistivity of $\rho=\dfrac{5\times10^{12}}{8.85}\ \Omega$m. The area of the plates is $A=1\text{m}^2$, and the dielectric film thickness $d=8.85\ \mu$m. The capacitor is charged from a 24 V supply connected in series with a 5 M$\Omega$ resistor. Dielectric constant of the material is 5.
\begin{qoptions}
\opt{A}{The maximum charge on the capacitor is 60$\mu$C.}
\opt{B}{The maximum charge on the capacitor is 120$\mu$C.}
\opt{C}{If the capacitor is disconnected from the circuit, the time in which the charge on capacitor is halved is 25 sec.}
\opt{D}{If the capacitor is disconnected from the circuit, the time in which the charge on capacitor is halved is 50 sec.}
\end{qoptions}
\qn{2} An electric field line is generated at charge $q_1$ and terminates at $q_2$ as shown. The equation of the line can be
\qfig{p2_phy_s1_q02_1}
\begin{qoptions}
\opt{A}{$q_1\left(1-\dfrac{x}{\sqrt{x^2+y^2}}\right)+q_2\left(1-\dfrac{\ell-x}{\sqrt{(\ell-x)^2+y^2}}\right)$\\$=q_1(1-\cos\alpha)$}
\opt{B}{$q_1\left(1-\dfrac{x}{\sqrt{x^2+y^2}}\right)-q_2\left(1-\dfrac{\ell-x}{\sqrt{(\ell-x)^2+y^2}}\right)$\\$=q_1(1-\cos\alpha)$}
\opt{C}{$q_1\left(1+\dfrac{x}{\sqrt{x^2+y^2}}\right)+q_2\left(1-\dfrac{\ell-x}{\sqrt{(\ell-x)^2+y^2}}\right)$\\$=q_1(1-\cos\alpha)$}
\opt{D}{$q_1\left(1+\dfrac{x}{\sqrt{x^2+y^2}}\right)-q_2\left(1-\dfrac{\ell-x}{\sqrt{(\ell-x)^2+y^2}}\right)$\\$=q_1(1-\cos\alpha)$}
\end{qoptions}
\qn{3} With what angular velocity the earth should spin in order that a body lying at $30^\circ$ latitude may become weightless [$R$ is radius of earth and $g$ is acceleration due to gravity on the surface of earth]
\begin{qoptions}
\opt{A}{$\sqrt{\dfrac{2g}{3R}}$}
\opt{B}{$\sqrt{\dfrac{4g}{R}}$}
\opt{C}{$\sqrt{\dfrac{g}{3R}}$}
\opt{D}{$\sqrt{\dfrac{4g}{3R}}$}
\end{qoptions}
\qn{4} A uniformly charged ring of charge $Q$ and radius $R$ is folded from its diameter such that two halves make an angle $60^\circ$ with each other then net electric field at centre of ring (O) is equal to
\qfig{p2_phy_s1_q04_1}
\begin{qoptions}
\opt{A}{$\dfrac{\sqrt3kQ}{\pi R^2}$}
\opt{B}{$\dfrac{2\sqrt3kQ}{\pi R^2}$}
\opt{C}{$\dfrac{2kQ}{\pi R^2}$}
\opt{D}{$\dfrac{4kQ}{\pi R^2}$}
\end{qoptions}
\sect{SECTION-I (ii)}
\qn{5} Circuits with non-linear components are of great importance. The circuit is designed to prevent damage to the device represented by the 10 k$\Omega$ load resistor when the voltage across the supply exceeds 2000 V. The component S has a V-I characteristic given by $I=kV^2$, where $V$ is the voltage across S when current $I$ flows through it. The constant $k$ has a value $k=1.0\times10^{-7}$ A V$^{-2}$.
\qfig{p2_phy_s2_q05_1}
\begin{qoptions}
\opt{A}{The current through the load resistor when operating normally under a supply of 2000 V is 0.1 A.}
\opt{B}{The current through the suppressor when operating normally under a supply of 2000 V is 0.2 A.}
\opt{C}{The supply voltage at which current through the load resistor doubles from its normal operating value is 3000 V.}
\opt{D}{The supply voltage at which current through the load resistor doubles from its normal operating value is 5000 V.}
\end{qoptions}
\qn{6} There is a solid metallic sphere of radius $R$, which is cut into two identical hemispheres. The cut surface is coated with a thin insulating layer of thickness $d$, and the two parts are put together so that the original shape of the sphere is restored. Initially, the sphere is electrically neutral. Then one of the hemispheres is given a positive charge $+Q$ while the other one remains neutral. You can assume $d\ll R$ in this problem.
\qfig{p2_phy_s2_q06_1}
\begin{qoptions}
\opt{A}{The surface charge density on the curved surface of the hemisphere is $\dfrac{Q}{4\pi R^2}$.}
\opt{B}{The electrostatic force on attraction between the two hemispheres is $\dfrac{3Q^2}{32\epsilon_0\pi R^2}$.}
\opt{C}{The electrostatic force of attraction between the two hemispheres is $\dfrac{3Q^2}{16\epsilon_0\pi R^2}$.}
\opt{D}{The electrostatic energy of the sphere is $\dfrac{Q^2}{2\pi\epsilon_0R}$}
\end{qoptions}
\qn{7} Three point masses having charge $q=3\times10^{-7}$ C and $m=9\times10^{-4}$ kg are connected to three strings each of length 20 cm such that the other end of the strings is connected to a common point on the ceiling as shown in the diagram. Initially string make an angle $\alpha=30^\circ$ with the vertical. System is released from rest.
\qfig{p2_phy_s2_q07_1}
\begin{qoptions}
\opt{A}{The angle made by the strings when the speed of point objects is maximum is $60^\circ$.}
\opt{B}{The maximum speed attained by the point objects is $\sqrt{2\left(\sqrt3+1\right)}$ m/s.}
\opt{C}{The maximum angle made by the strings with the vertical is $90^\circ$.}
\opt{D}{The maximum tangential acceleration of point objects is 40 m/s$^2$.}
\end{qoptions}
\qn{8} A potentiometer wire PQ is made by using two wires of different cross sectional area $A$ and $2A$. The material of wire PS and SQ is same and their lengths are 5m each. A source battery of 10V and rheostat $R_h$ are connected in the circuit as shown in the figure. Choose the correct option(s) regarding the balance point.
\qfig{p2_phy_s2_q08_1}
\begin{qoptions}
\opt{A}{If a cell of emf $e_0=5$V is used then balance point from the end P is at 7.5 m}
\opt{B}{If a cell of emf $e_0=5$V is used then balance point from the end P is at 3.75 m}
\opt{C}{If a cell of emf $e_0=8$V is used then balance point from the end P is at 7 m}
\opt{D}{If a cell of emf $e_0=8$V is used then balance point from the end P is at 8.5 m}
\end{qoptions}
\qn{9} Initially all the switch in the circuit are open.
\qfig{p2_phy_s2_q09_1}
\begin{qoptions}
\opt{A}{The current in AB after closing $K_1$ is $\dfrac{V_0}{2R}$.}
\opt{B}{The current in AC after closing $K_1$ and $K_2$ is $\dfrac{V_0}{2R}$.}
\opt{C}{The current in BC after closing all the switches is $\dfrac{V_0}{2R}$.}
\opt{D}{The current in AB after closing all the switches is $\dfrac{V_0}{R}$.}
\end{qoptions}
\qn{10} A conductor is made of an isotropic material (resistivity $\rho$) has rectangular cross-section. Horizontal dimension of the rectangle decreases linearly from $2x$ at one end to $x$ at the other end and vertical dimension increases from $y$ to $2y$ as shown in figure. Length of the conductor along the axis is equal to $\lambda$. A battery is connected across the ends of this conductor then
\qfig{p2_phy_s2_q10_1}
\begin{qoptions}
\opt{A}{rate of generation of heat per unit length is maximum at the ends of conductor}
\opt{B}{rate of generation of heat per unit length is maximum at middle cross-section}
\opt{C}{drift velocity of conduction electrons is minimum at middle section}
\opt{D}{electric field intensity is maximum at ends of conductor.}
\end{qoptions}
\sect{SECTION-III}
\qn{1} For the circuit shown in figure, what should be the value of the variable resistor $R$ (in $\Omega$) in order to minimise the power in the $10\,\Omega$ resistor. Fill the value of $\dfrac{R}{20}$.
\qfig{p2_phy_s3_q01_1}
\qn{2} A charged particle having unity specific charge enters a system of three parallel meshes. The dimensions of meshes are very large as compared to the gap between the meshes. The meshes are connected to batteries of emf $E_1=15$ V and $E_2=12$ V as shown in the diagram.
\qfig{p2_phy_s3_q02_1}
The charged particle comes to rest at a point $\dfrac d4$ behind the second mesh. The speed of particle at a large distance from the meshes is.
\qn{3} The key K is kept open for a very long time. The current through the capacitance $C$ just after the key K is closed is $\dfrac{x}{10}$A. The value of $x$ is. Take : $E_1=40$ V, $r_1=200\ \Omega$, $E_2=80$ V, $r_2=50\Omega$ and $R=150\ \Omega$.
\qfig{p2_phy_s3_q03_1}
\qn{4} The time period of Jupiter is $(125)^{1/2}$ years. The distance of Jupiter from Sun in terms of astronomical unit is given by $\alpha$ AU. Find the value of $\alpha$. (Astronomical unit is defined as the average distance of Earth from the Sun).
\qn{5}
\qfig{p2_phy_s3_q05_1}
Three concentric conducting spheres A (solid sphere of radius $R$), B (hollow sphere of inner radius $2R$ and outer radius $3R$) and C (hollow sphere of radius $4R$) are kept as shown. Sphere A is given charge $Q$. Capacity of the system is as follows:\\
$C_1$ : When both $S_1$ and $S_2$ are open\\
$C_2$ : When $S_1$ is closed and $S_2$ is open\\
$C_3$ : When $S_1$ is open and $S_2$ is closed\\
$C_4$ : When both $S_1$ and $S_2$ are close\\
If $(C_3+C_4)-(C_1+C_2)=\dfrac{\alpha\beta\pi\varepsilon_0R}{15}$ (Where $\alpha\beta$ is a two digit number) then calculate $\beta-\alpha=$ \underline{\hspace{1.5cm}}.
\qn{6} A narrow tunnel is dug inside the earth along one of its diameters. Radius of earth is $R$ and its mass is $M$. A particle is projected inside the tunnel with velocity $\sqrt{\dfrac{2GM}{R}}$ from one of its ends then maximum velocity attained by the particle in the subsequent motion is $\sqrt{\dfrac{nGM}{R}}$ (assuming tunnel to be frictionless). Find $n$.
\qn{7} The potential difference $V_B-V_A$ (in V) for the circuit shown in the figure is
\qfig{p2_phy_s3_q07_1}
\qn{8} Two conducting spheres of radii $r$ and $2r$ are placed at very large separation. Each sphere possesses charge $Q$. These spheres are connected with a conducting wire of resistance $R$ then current reduces to half the initial current after time $\dfrac{8\pi\epsilon_0 rR\,\ell n2}{n}$. Find $n$.
\partsec{2}{CHEMISTRY}
\sect{SECTION-I (i)}
\qn{1} Consider the four different galvanic cells having equal cell potential in given conditions :\\
Ag $|$ Ag$^+$ (Ag$_2$CrO$_4$, saturated sol.) $\|$ Ag$^+$ (0.1M) $|$ Ag, $E_{cell}=xV$ and $K_{sp,Ag_2CrO_4}=$`p'\\
Ag $|$ Ag$^+$(AgCl saturated sol.) $\|$ Ag$^+$(0.1M) $|$ Ag, $E_{cell}=xV$ and $K_{sp,AgCl}=$`q'\\
Ag $|$ Ag$^+$(Ag$_3$PO$_4$ saturated sol.) $\|$ Ag$^+$(0.1M) $|$ Ag, $E_{cell}=xV$ and $K_{sp,Ag_3PO_4}=$`r'\\
Ag $|$ Ag$^+$ (1M) $\|$ Ag$^+$ (Ag$_2$SO$_4$ saturated sol.) $|$ Ag, $E_{cell}=xV$ and $K_{sp,Ag_2SO_4}=$`s'\\
The correct option for solubility product for all the four cases is (by considering all cells are in equilibrium stage)
\begin{qoptions}
\opt{A}{$s>q>p>r$}
\opt{B}{$s<p<q<r$}
\opt{C}{$r>s>q>p$}
\opt{D}{$r>q>p>s$}
\end{qoptions}
\qn{2} For an electrolytic cell undergoing following balanced reaction :\\
NO$_3^-$ + 3H$_3$O$^+$ + Hg$_2^{2+}\to$ 2Hg$^{2+}$ + HNO$_2$ + 4H$_2$O\\
If an average current of 10 amp. passes through the cell, how long will it take to produce 0.3 moles of HNO$_2$ by this reaction?
\begin{qoptions}
\opt{A}{1.6 hours}
\opt{B}{16 hours}
\opt{C}{3.2 hours}
\opt{D}{32 hours}
\end{qoptions}
\qn{3} For the reaction $3A\to2B$, the following graph is obtained experimentally. The rate of reaction when $[A]=0.2$ M, is
\qfig{p2_chem_s1_q03_1}
\begin{qoptions}
\opt{A}{$8\times10^{-4}$ M min$^{-1}$}
\opt{B}{$\dfrac{16}{6}\times10^{-4}$ M min$^{-1}$}
\opt{C}{$\dfrac43\times10^{-3}$ M min$^{-1}$}
\opt{D}{$\dfrac{16}{3}\times10^{-3}$ M min$^{-1}$}
\end{qoptions}
\qn{4} Which of the following is colloid in which dispersed phase is gas and dispersion medium is liquid ?
\begin{qoptions}
\opt{A}{Mist}
\opt{B}{Ruby Glass}
\opt{C}{Whipped Cream}
\opt{D}{Lava}
\end{qoptions}
\sect{SECTION-I (ii)}
\qn{5} NH$_3$ concentration is a saturated solution of AgBr(s) is maintained at 0.2M. The correct information(s) related with the solution is/are :\\
Consider the data given below\\
$K_{sp}$ of AgBr $=2\times10^{-10}$\\
$K_f$ of Ag(NH$_3$)$_2^+=8\times10^6$\\
$\lambda^o_{(Ag^+)}=6\times10^{-3}$ Sm$^2$ mol$^{-1}$\\
$\lambda^o_{(Br^-)}=8\times10^{-3}$ Sm$^2$ mol$^{-1}$
\begin{qoptions}
\opt{A}{Conductivity of Br$^-=6.4\times10^{-2}$ S m$^{-1}$}
\opt{B}{Conductivity of Ag$^+=1.5\times10^{-7}$ S m$^{-1}$}
\opt{C}{$[\mathrm{Ag^+}]=8\times10^{-3}$M}
\opt{D}{$[\mathrm{Br^-}]=2\times10^{-3}$M}
\end{qoptions}
\qn{6} EMF of the cell: Cd(s) $|$ CdCl$_2$.5H$_2$O (saturated solution) $|$ AgCl(s) $|$ Ag(s) is $+0.70$ V at $0^\circ$C and $+0.60$ V at $50^\circ$C. If $\Delta H^\circ$ and $\Delta S^\circ$ are temperature independent, then the correct information(s) regarding the cell reaction is/are
\begin{qoptions}
\opt{A}{$\Delta G^\circ=-115.8$ kJ at $50^\circ$C}
\opt{B}{$\Delta G^\circ=135.1$ kJ at $0^\circ$C}
\opt{C}{$\Delta S^\circ=-386$ J/K}
\opt{D}{$\Delta H^\circ=221.18$ kJ}
\end{qoptions}
\qn{7} A substance A reacts at 500 K by parallel reactions to give B and C as shown
\qfig{p2_chem_s2_q07_1}
\hspace*{1.2em}\textbf{Reaction -1} : $k_1=10^{-3}$ sec$^{-1}$. $E_{A_1}=80$ kcal/mole\\
\hspace*{1.2em}\textbf{Reaction -2} : $k_2=10^{-2}$ sec$^{-1}$. $E_{A_2}=100$ kcal/mole\\
Identify the \textbf{INCORRECT} statement(s).
\begin{qoptions}
\opt{A}{The net activation energy for conversion of \textbf{A} is 90 Kcal/mole.}
\opt{B}{The pre-exponential factor for 1$^{st}$ reaction is smaller as compared to 2$^{nd}$ reaction.}
\opt{C}{$|\Delta H|$ for \textbf{reaction-1} must be greater than \textbf{reaction$-$2}.}
\opt{D}{Fraction of \textbf{A} reacting by either of the reaction will be same at 500 K.}
\end{qoptions}
\qn{8} For the reaction CHCl$_3$ + Cl$_2\xrightarrow{h\nu}$ CCl$_4$ + HCl, mechanism is given below :\\
Cl$_{2(g)}\ \xrightleftharpoons[K_{-1}]{K_1}\ 2$Cl$_{\mathrm{(g)}}$\\
Cl$_{\mathrm{(g)}}$ + CHCl$_{3\mathrm{(g)}}$ $\xrightarrow{k_2}$ HCl$_{\mathrm{(g)}}$ + CCl$_{3\mathrm{(g)}}$\\
CCl$_{3\mathrm{(g)}}$ + Cl$_{\mathrm{(g)}}$ $\xrightarrow{k_3}$ CCl$_{4\mathrm{(g)}}$\\
$k_1=4.8\times10^3$ sec$^{-1}$, $K_{-1}=12\times10^3$ M$^{-1}$ sec$^{-1}$, $k_2=1.3\times10^{-2}$ M$^{-1}$ sec$^{-1}$, $k_3=2.1\times10^2$ M$^{-1}$ sec$^{-1}$\\
Given that second step of the mechanism is slowest step.\\
Correct statement is/are :
\begin{qoptions}
\opt{A}{Order of reaction is 3/2}
\opt{B}{Magnitude of overall rate constant (conc. in M, time in sec) is $2.6\times10^{-2}$}
\opt{C}{If concentration of CHCl$_3$ is increased to four times rate of reaction increases by a factor of two}
\opt{D}{If concentration of Cl$_2$ is increased to four time rate of reaction increase by a factor of two}
\end{qoptions}
\qn{9} Which of the following statement(s) is/are incorrect with respect to enzyme catalysed reactions?
\begin{qoptions}
\opt{A}{The cavities on enzyme ``Maltase'' \& structure of molecules of Maltose have complimentary shapes for conversion of Maltose to glucose.}
\opt{B}{Just like other catalysts, catalytic action of enzyme is more at a lower temperature.}
\opt{C}{In presence of vitamins \& some metal ions; enzyme activity generally increases.}
\opt{D}{Most of the enzyme catalysed reactions are elementary.}
\end{qoptions}
\qn{10} For the coagulation of a particular colloidal solution the order of coagulation value of electrolytes is as follows.
\[\mathrm{Na_3PO_4>BaSO_4>AlCl_3}\]
\begin{qoptions}
\opt{A}{In electroosmosis of this colloidal solution dispersion medium move towards cathode.}
\opt{B}{Diffused layer of the colloidal sol particles contains positively charged ions.}
\opt{C}{It may be a metal sol such as Pt, Au or Hg sol}
\opt{D}{It may be basic dye sol such as methylene blue.}
\end{qoptions}
\sect{SECTION-III}
\qn{1} Consider the strong electrolytes Z$_m$X$_n$, U$_m$Y$_p$ and V$_m$X$_n$ . Limiting molar conductivity ($\Lambda^0$) of Z$_m$X$_n$ and V$_m$X$_n$ are 250 and 440 S cm$^2$ mol$^{-1}$, respectively. The nearest integer value of $(m+n+p)$ is \underline{\hspace{1.2cm}}.\\
Given :
\par\nobreak
{\small\setlength{\tabcolsep}{3pt}\renewcommand{\arraystretch}{1.3}
\begin{tabular}{|L{2.4cm}|c|c|c|c|c|}\hline
Ion & Z$^{n+}$ & U$^{p+}$ & V$^{n+}$ & X$^{m-}$ & Y$^{m-}$\\\hline
$\lambda^0$ (S cm$^2$ mol$^{-1}$) & 50.0 & 25.0 & 100.0 & 80.0 & 100.0\\\hline
\end{tabular}}\\
$\lambda^{\,0}$ is the limiting molar conductivity of ions\\
The plot of molar conductivity ($\Lambda$) of U$_m$Y$_p$ $vs$ $c^{1/2}$ is given below.
\qfig{p2_chem_s3_q01_2}
\qn{2} For the reaction : $2$NO(g) + H$_2$(g) $\to$ N$_2$O(g) + H$_2$O(g), the value of $\dfrac{-dp}{dt}$ was found to be 1.5 PaS$^{-1}$ for a pressure of 372 Pa for NO and 0.25 PaS$^{-1}$ for a pressure of 152 Pa for NO, the pressure of H$_2$ being constant. If pressure of NO was kept costant, the value of $\dfrac{-dp}{dt}$ was found 1.60 PaS$^{-1}$ for a pressure of 289 Pa for H$_2$ and 0.79PaS$^{-1}$ for a pressure of 144 Pa for H$_2$. If order of reaction with respect to NO and H$_2$ are $a$ and $b$, respectively, then find nearest integer value of $(a+b)$ ?
\qn{3} $K_{eq}$ for the reaction :\\
A(s) + 3B$^{+}$(aq.) $\rightleftharpoons$ A$^{3+}$(aq.) + 3B(s) is $10^x$ at 25\textdegree C. `x' is\\
Given : $\dfrac{2.303RT}{F}=0.06$V , $E^0_{A/A^{2+}}=0.25$V , $E^0_{A^{2+}/A^{3+}}=0.1$V , $E^0_{B/B^+}=0.1$V
\qn{4} The voltage of the cell is $-0.6769$ V at 25\textdegree C temperature.\\
Pt(s)$|$H$_2$(g) $|$ NaHSO$_3$(aq), Na$_2$SO$_3$ $\|$ Zn$^{2+}$(aq) $|$ Zn(s)\\
\hspace*{0.6em}1 atm \hspace{1.2em} 0.4M \hspace{0.6em} $6.44\times10^{-3}$M \hspace{0.4em} 0.04M\\
Also Zn$^{2+}$ + 2e$^-$ $\to$ Zn(s), $E^\circ=-0.736$ V\\
If $K=\dfrac{[\mathrm{H^+}]\,[\mathrm{SO_3^{2-}}]}{[\mathrm{HSO_3^-}]}=3.22\times10^{-x}$. Then find value of `x' ?\\
$\left[\text{Given : }\dfrac{2.303RT}{F}=0.0591\text{ V}\right]$
\qn{5} CH$_3$COOC$_2$H$_{5\mathrm{(aq)}}$ + H$_2$O$_{\mathrm{(\ell)}}$ $\xrightarrow{H^+}$ CH$_3$COOH$_{(aq)}$ + C$_2$H$_5$OH$_{(aq)}$\\
Rate $=K[\mathrm{H^+}]^1[\mathrm{H_2O}]^1[\mathrm{CH_3COOC_2H_5}]^1$, $K=1.8\times10^{-3}$ M$^{-2}$ sec$^{-1}$.\\
In an experiment 1 L of 1 M CH$_3$COOC$_2$H$_5$ is hydrolysed by mixing 1 L 0.2 M HCl. Immediately after starting the reaction, $V$ mL of reaction mixture requires 10 mL 0.1 M NaOH, after long time of starting the reaction $V$ mL of reaction mixture requires 100 ml of 0.1 M NaOH for complete neutralisation. If volume of 0.1 M NaOH required for complete neutralisation of $V$ mL of reaction mixture after 230.3 sec is `$x$' mL find nearest integer value of $\left(\dfrac{x}{10}\right)$?
\qn{6} Total number of sols among below given options which can be coagulated by Al$^{+3}$ ion are?\\
Eosin, methylene blue, Congo red, TiO$_2$, Ag Sol, Sb$_2$S$_3$ Sol, Al$_2$O$_3$ sol. Haemoglobin sol, Clay sol., Charcoal sol.
\qn{7} For a reaction\\
A(g) $\xrightarrow{\text{Rate constant}=k}$ P(g) + Q(g)\\
carried in closed and rigid container at constant temperature $T$, following informations are provided.
\par\nobreak
{\small\setlength{\tabcolsep}{4pt}\renewcommand{\arraystretch}{1.3}
\begin{tabular}{|c|c|c|c|}\hline
 & Initial Concentration & $T_{1/2}$ & $T_{3/4}$\\\hline
A(g) & $1\times10^{-2}$M & 10 Min & 30 Min\\\hline
\end{tabular}}\\
If initial rate of reaction in mole/lit/min is represented in scientific notation as $x\times10^{-y}$ then find the value of $y$.
\qn{8} 1 gm of charcoal having surface area $3.01\times10^2$ m$^2$ is mixed in 100 ml of 0.5 M acetic acid solution; after sufficient time concentration of solution becomes 0.4 M. The surface area occupied by one molecule of acetic acid is $y\times10^{-20}$ m$^2$. Then $y$ is. [Given : $N_A=6.02\times10^{23}$]
\partsec{3}{MATHEMATICS}
\sect{SECTION-I (i)}
\qn{1} Let $A_1$, $A_2$ and $A_3$ be the regions on $\mathbb R^2$ defined by\\
$A_1=\{(x,y):x\ge0,\ y\ge0,\ 2x+2y-x^2-y^2>1>x+y\}$,\\
$A_2=\{(x,y):x\ge0,\ y\ge0,\ x+y>1>x^2+y^2\}$,\\
$A_3=\{(x,y):x\ge0,\ y\ge0,\ x+y>1>x^3+y^3\}$\\
Denote by $|A_1|,|A_2|$ and $|A_3|$ the areas of the regions $A_1$, $A_2$ and $A_3$ respectively. Then
\begin{qoptions}
\opt{A}{$|A_1|>|A_2|>|A_3|$}
\opt{B}{$|A_1|>|A_3|>|A_2|$}
\opt{C}{$|A_1|=|A_2|<|A_3|$}
\opt{D}{$|A_1|=|A_3|>|A_2|$}
\end{qoptions}
\qn{2} Solution of the differential equation $\dfrac{\sqrt x\,dx+\sqrt y\,dy}{\sqrt x\,dx-\sqrt y\,dy}=\sqrt{\dfrac{y^3}{x^3}}$ is given by (Base of all $\log=e$)
\begin{qoptions}
\opt{A}{$\dfrac32\log\left(\dfrac yx\right)+\log\left|\dfrac{x^{3/2}+y^{3/2}}{x^{3/2}}\right|+\tan^{-1}\left(\dfrac yx\right)^{3/2}+c=0$}
\opt{B}{$\dfrac23\log\left(\dfrac yx\right)+\log\left|\dfrac{x^{3/2}+y^{3/2}}{x^{3/2}}\right|+\tan^{-1}\left(\dfrac yx\right)^{3/2}+c=0$}
\opt{C}{$\dfrac23\log\left(\dfrac yx\right)+\log\left(\dfrac{x+y}{x}\right)+\tan^{-1}\left(\dfrac{y^{3/2}}{x^{3/2}}\right)+c=0$}
\opt{D}{$\dfrac12\log\left(x^3+y^3\right)+\tan^{-1}\left(\dfrac yx\right)^{3/2}+c=0$}
\end{qoptions}
\qn{3} Let $I_1=\displaystyle\int_{-2}^{2}\dfrac{x^6+3x^5+7x^4}{x^4+2}\,dx$ and $I_2=\displaystyle\int_{-3}^{1}\dfrac{2(x+1)^2+11(x+1)+14}{(x+1)^4+2}\,dx$, then the value of $I_1+I_2$ is
\begin{qoptions}
\opt{A}{8}
\opt{B}{200/3}
\opt{C}{100/3}
\opt{D}{None of these}
\end{qoptions}
\qn{4} $\displaystyle\int\dfrac{x\left(x\tan^{-1}x+(\ell nx)(\ell n(\ell nx))\right)+\tan^{-1}x}{(x^3+x)\ell nx}\,dx=\phi(x)$ where $\phi(e)=0$ then the value of $\displaystyle\lim_{x\to1}\dfrac{\phi(x)}{\tan\left(\frac{\pi x}{2}\right)}$ is equal to
\begin{qoptions}
\opt{A}{0}
\opt{B}{$\dfrac\pi4$}
\opt{C}{$-\dfrac\pi4$}
\opt{D}{$-\dfrac12$}
\end{qoptions}
\sect{SECTION-I (ii)}
\qn{5} For $0<a<1$, $-1<x<1$. Let $f(x)=\dfrac{a-x}{1-ax}$ and if $\displaystyle\int_0^a\dfrac{1-(f(x))^6}{1-x^2}\,dx=F(a)$ then
\begin{qoptions}
\opt{A}{$F(a)$ is increasing function}
\opt{B}{$F(a)$ is an odd function}
\opt{C}{$F(1)=\dfrac{23}{15}$}
\opt{D}{$F(x)=0$ has exactly one real root}
\end{qoptions}
\qn{6} Let $y=\sin x$ and $I=\displaystyle\int(x-y)\tan^2\left(\dfrac x2\right).d(x+y)$, then
\begin{qoptions}
\opt{A}{$d(y+x)=\cot^2\left(\dfrac x2\right).d(x-y)$}
\opt{B}{$\dfrac{d^2x}{dy^2}=\dfrac{\sin x}{\cos^3x}$}
\opt{C}{$I=(x-y)^2+C$}
\opt{D}{$I=\dfrac{(x-y)^2}{2}+C$}
\end{qoptions}
\qn{7} The area bounded by the function $f(x)=|\cos x|$ and $g(x)=\sin x$ between the ordinates $x=0$ and $x=n\pi$, $n\in N$, is equal to
\begin{qoptions}
\opt{A}{$4n\sqrt2$ if $n$ is odd}
\opt{B}{$2n\sqrt2$ if $n$ is even}
\opt{C}{$2\sqrt2\,(n+1)$ if $n$ is odd}
\opt{D}{$2\sqrt2\,(n+1)-4$ if $n$ is odd}
\end{qoptions}
\qn{8} If the Area bounded by $y=(\sin x)^{\mathrm{cosec}\,x}$, $y=(\mathrm{cosec}\,x)^{\sin x}$, $y=(\sin x)^{\sin x}$ and $y=(\mathrm{cosec}\,x)^{\mathrm{cosec}\,x}$ between the ordinates $x=0$, $x=\dfrac\pi2$ denoted by $A_1$, $A_2$, $A_3$ and $A_4$ then
\begin{qoptions}
\opt{A}{$A_4$ is the greatest}
\opt{B}{$A_1$ is the least}
\opt{C}{$A_2>A_3$}
\opt{D}{$A_2<A_3$}
\end{qoptions}
\qn{9} Let $I_n=\displaystyle\int\dfrac{\cos^{2n}x}{\sin x}\,dx$ satisfies the recursive relation $(2n+1)I_{n+1}=(2n+1)I_n+\phi(x)$, if $\phi'(x)=f(n,x)$, then which is/are correct
\begin{qoptions}
\opt{A}{$f\left(5,\dfrac\pi2\right)=0$}
\opt{B}{$f\left(5,\dfrac\pi4\right)=\dfrac{-11}{32\sqrt2}$}
\opt{C}{$f(6,\pi)$ is even number}
\opt{D}{$f\left(6,\dfrac\pi2\right)$ is prime number}
\end{qoptions}
\qn{10} Which is/are correct ?
\begin{qoptions}
\opt{A}{$\displaystyle\int_0^2\dfrac{1}{2+\arctan x}\,dx\ge\ell n2$}
\opt{B}{$\displaystyle\int_0^2\dfrac{1}{2+\arctan x}\,dx<\ell n2$}
\opt{C}{If $f(x)$ is such that $|f(u)-f(v)|\le|u-v|$ then $\left|\displaystyle\int_a^b f(x)\,dx-(b-a).f(a)\right|\le\dfrac{(b-a)^2}{2}$, $(a<b)$}
\opt{D}{If $f(x)$ is such that $|f(u)-f(v)|\le|u-v|$ then $\left|\displaystyle\int_a^b f(x)\,dx-(b-a)f(a)\right|\ge\dfrac{(b-a)^2}{2}$, $(a<b)$}
\end{qoptions}
\sect{SECTION-III}
\qn{1} The value of $\displaystyle\int_{\pi/4}^{65\pi/4}\dfrac{dx}{(1+2^{\cos x})(1+2^{\sin x})}$ is $k\pi$, then the value of $k$ is equal to
\qn{2} Let for $x\in R$, $S_0(x)=x$, $S_k(x)=C_kx+k\displaystyle\int_0^x S_{k-1}(t)\,dt$, where $C_0=1$, $C_k=1-\displaystyle\int_0^1 S_{k-1}(x)\,dx$, $k=1,2,3,\ldots$ Then $\left[\dfrac{S_2(3)+6C_3}{3}\right]$ is equal to (Where $[\ ]$ denotes Greatest Integer Function)
\qn{3} If $\displaystyle\int\dfrac{3x^2+2x}{x^6+2x^5+x^4+2x^3+2x^2+5}\,dx=F(x)$, then find the value of $[F(1)-F(0)]$, where $[\,.\,]$ represents greatest integer function
\qn{4} Let the functions $f:R\to R$ and $g:R\to R$ be defined by $f(x)=e^{x-1}-e^{-|x-1|}$ and $g(x)=\dfrac12\left(e^{x-1}+e^{1-x}\right)$. Then the area of the region bounded by the curves $y=f(x)$, $y=g(x)$ and $x=1$ is $2-\sqrt\lambda$, then the value of $\lambda$ is
\qn{5} Let $a$, $b$ and $\lambda$ be positive real numbers. Suppose $P$ is an end point of the latus rectum of the parabola $y^2=4\lambda x$, and suppose the ellipse $\dfrac{x^2}{a^2}+\dfrac{y^2}{b^2}=1$ passes through the point $P$. If the tangents to the parabola and the ellipse at the point $P$ are perpendicular to each other, then the eccentricity of the ellipse is $e$. If $\left(e,\dfrac ab\right)$ is a point on the ellipse $15x^2+3y^2=K$, then $K$ is equal to \underline{\hspace{1.5cm}}
\qn{6} Let $\dfrac{dy}{dx}=\dfrac{ax-by+a}{bx+cy+a}$, $a,b,c\in R$ represents a circle with centre $(\alpha,\beta)$ then $\alpha+2\beta=$
\qn{7} The area bounded by the curve $y=f(x)=\displaystyle\int_1^{1/x^3}\left(\dfrac1{1+t^4}+t^{-4/3}\right)dt-3\int_x^1\dfrac{t^8}{1+t^{12}}\,dt$ and the co-ordinate axes is equal to $\dfrac pq$ (where $p,q\in N$, G.C.D $=1$) then the value of $p+q$ is
\qn{8} Let $u=\displaystyle\int_0^{\pi/4}\left(\dfrac{\cos x}{\sin x+\cos x}\right)^2dx$ and $v=\displaystyle\int_0^{\pi/4}\left(\dfrac{\sin x+\cos x}{\cos x}\right)^2dx$, then the value of $\dfrac vu$ equals
\end{multicols}
\end{document}
